\documentclass[11pt,a4paper]{article}

\usepackage[a4paper, margin=2cm]{geometry}
\usepackage{graphicx}
\usepackage{booktabs}
\usepackage{hyperref}
\usepackage{float}
\usepackage{amsmath}
\usepackage{amssymb}
\usepackage{xcolor}
\usepackage{listings}
\usepackage{caption}
\usepackage{subcaption}
\usepackage{parskip}
\usepackage{enumitem}
\usepackage{titlesec}
\usepackage{fancyhdr}
\usepackage{microtype}

% Header/Footer
\pagestyle{fancy}
\fancyhf{}
\rhead{SMAI Assignment 3}
\lhead{Bengali Book Review Sentiment}
\cfoot{\thepage}

% Hyperlink styling
\hypersetup{
    colorlinks=true,
    linkcolor=blue!70!black,
    urlcolor=blue!70!black,
    citecolor=blue!70!black
}

% Code listing style
\lstset{
    basicstyle=\ttfamily\small,
    backgroundcolor=\color{gray!10},
    frame=single,
    breaklines=true,
    keywordstyle=\color{blue},
    commentstyle=\color{green!60!black},
    stringstyle=\color{orange!80!black},
    language=Python
}

\title{
    \textbf{Bengali Book Review Sentiment Analysis} \\[0.5em]
    \large Fine-tuning IndicBERT for Low-Resource Indic NLP \\[0.3em]
    \normalsize SMAI Assignment 3 --- T8.4
}
\author{Souradeep Das (2025201004) \\
\texttt{souradeep.das@students.iiit.ac.in} \\
Kushal Mukherjee (2025201072) \\
\texttt{kushal.mukherjee@students.iiit.ac.in} \\
Monalisa Sarkar (2024801018) \\
\texttt{monalisa.sarkar@research.iiit.ac.in} \\
\textit{Team: Axiom AI}}
\date{May 5, 2026}

\begin{document}

\maketitle
\thispagestyle{fancy}

\noindent Project repository: \url{https://github.com/Kushal-2002/Bengali-Sentiment-Analysis} \\
Model on Hugging Face: \url{https://huggingface.co/k040902/bengali_book_review_sentiment} \\
Live app: \url{https://smai-assignment-3.streamlit.app/}

%-----------------------------------------------------------
\section{Introduction}
%-----------------------------------------------------------

Sentiment analysis — the task of automatically classifying the opinion expressed in a piece of text as positive, negative, or neutral — is one of the most practical applications of Natural Language Processing (NLP). For readers, authors, and publishers, being able to automatically summarize large volumes of Bengali book reviews helps track reader reception, identify recurring themes, and surface strengths or weaknesses across titles.

While sentiment analysis is well-studied for English, the problem is significantly harder for Indic languages such as Bengali (\textit{Bangla}). Bengali is the seventh most spoken language in the world with over 230 million speakers, yet NLP resources and pre-trained models for it remain limited compared to English. Challenges include:
\begin{itemize}[noitemsep]
    \item \textbf{Script complexity:} Bengali uses the Brahmic script with complex ligatures and conjunct characters.
    \item \textbf{Morphological richness:} Words can have many inflected forms, making vocabulary coverage harder.
    \item \textbf{Limited labelled data:} Publicly available, high-quality labelled Bengali sentiment datasets are scarce.
    \item \textbf{Code-mixing:} Real reviews often mix Bengali with English words, adding further difficulty.
\end{itemize}

This report describes the design, implementation, and evaluation of a Bengali sentiment analyser for book reviews that fine-tunes the multilingual IndicBERT model on Bengali review data, and exposes it through an interactive Streamlit web application.

%-----------------------------------------------------------
\section{Dataset}
%-----------------------------------------------------------

\subsection{Source}

We use the \textbf{AI4Bharat IndicSentiment} dataset~\cite{indicsentiment}, a curated collection of Bengali reviews annotated for sentiment. For this assignment, we use the Bengali (\texttt{bn}) subset.

\subsection{Split Usage}

The assignment instructions specify a non-standard split convention:
\begin{center}
\begin{tabular}{lll}
\toprule
\textbf{HuggingFace Split} & \textbf{Role in this work} & \textbf{Size} \\
\midrule
\texttt{test}       & Used for \textbf{training} the model & 998 samples \\
\texttt{validation} & Used as the held-out \textbf{test set} & 156 samples \\
\bottomrule
\end{tabular}
\end{center}

This reversal is deliberate: the larger \texttt{test} split is used for training to maximise the data available for fine-tuning, while the \texttt{validation} split is kept untouched as the final evaluation set.

From the 998-sample training pool we further carve a stratified \textbf{10\% validation split} (100 samples), leaving 898 samples for training. This inner split is what selects the best epoch; the 156-sample test set is touched exactly once, after training, to produce the numbers reported in Section~\ref{sec:results}. Keeping epoch selection off the test set is essential: an earlier version of this work selected the best checkpoint by scoring on the test set itself, which inflates the reported accuracy because the test set has then influenced model selection.

\subsection{Label Distribution}

The dataset contains binary sentiment labels: \textbf{Positive} and \textbf{Negative}. No neutral class is present in the Bengali subset. The training set has a slight imbalance towards positive reviews, which is typical of review datasets where satisfied readers are often more motivated to leave feedback.

\subsection{Preprocessing}

Each sample consists of a Bengali review string (column \texttt{INDIC REVIEW}) and a string label (column \texttt{LABEL}). Rows with missing values in either column are dropped. No additional text cleaning is applied, allowing the pre-trained tokenizer to handle tokenisation natively.

%-----------------------------------------------------------
\section{Method}
%-----------------------------------------------------------

\subsection{Model: IndicBERT}

We use \textbf{IndicBERT} (\texttt{ai4bharat/indic-bert}), a multilingual ALBERT-based model pre-trained by AI4Bharat on a large Indic language corpus covering 12 major Indian languages including Bengali. Key properties:

\begin{itemize}[noitemsep]
    \item \textbf{Architecture:} ALBERT (A Lite BERT) — uses parameter sharing across layers and factorised embedding to reduce model size while maintaining performance.
    \item \textbf{Pre-training corpus:} IndicCorp, a large web-crawled corpus of Indic text (~9 billion tokens).
    \item \textbf{Tokenizer:} SentencePiece with a shared vocabulary across all 12 languages, enabling sub-word tokenisation robust to morphological variation.
    \item \textbf{Why IndicBERT over mBERT?} IndicBERT is specifically optimised for Indic scripts and has dedicated vocabulary coverage for Bengali, making it substantially better than generic multilingual models for this task.
\end{itemize}

For classification, a linear head (\texttt{AlbertForSequenceClassification}) is added on top of the \texttt{[CLS]} token representation.

\subsection{Tokenisation}

Input texts are tokenised using the IndicBERT SentencePiece tokenizer with:
\begin{itemize}[noitemsep]
    \item \textbf{Truncation:} Sequences longer than 128 tokens are truncated. Most reviews are well within this limit.
    \item \textbf{Dynamic padding:} Batches are padded to the longest sequence in each batch (not a fixed global max), reducing unnecessary computation.
    \item \textbf{Label encoding:} String labels (\texttt{"Positive"}, \texttt{"Negative"}) are mapped to integer IDs ($0, 1$) via a deterministic \texttt{label2id} dictionary built from the sorted unique labels in the training set.
\end{itemize}

\subsection{Training Setup}

\begin{center}
\begin{tabular}{ll}
\toprule
\textbf{Hyperparameter} & \textbf{Value} \\
\midrule
Optimizer               & AdamW (default in HuggingFace Trainer) \\
Learning rate           & $2 \times 10^{-5}$ \\
Batch size (train/eval) & 16 \\
Epochs                  & 8 \\
Random seed             & 42 \\
Max sequence length     & 128 tokens \\
Evaluation strategy     & Each epoch; best by macro F1 on validation \\
Checkpoint saving       & Every epoch, keeping last 1 \\
\bottomrule
\end{tabular}
\end{center}

The learning rate of $2 \times 10^{-5}$ is the standard fine-tuning rate for BERT-family models. Training for 8 epochs on a dataset of $\sim$1000 samples is sufficient to converge without severe overfitting.

\subsection{Evaluation Metric}

We report \textbf{accuracy} (fraction of correctly classified samples) and \textbf{macro F1} (the unweighted mean of the per-class F1 scores) on the held-out test set. Accuracy is easy to interpret and appropriate here as the binary label distribution is not severely skewed; macro F1 is reported alongside it because it weights both classes equally and therefore exposes a model that does well on the majority class only. Macro F1 on the validation split is also the criterion used to select the best epoch during training.

%-----------------------------------------------------------
\section{Results}
\label{sec:results}
%-----------------------------------------------------------

\subsection{Training Dynamics}

Training ran for 8 epochs (456 optimiser steps, 272.8 seconds). The training loss fell steadily from 0.68 to 0.15, while validation accuracy and macro F1 improved over the same period:

\begin{center}
\begin{tabular}{ccccc}
\toprule
\textbf{Epoch} & \textbf{Train loss} & \textbf{Val loss} & \textbf{Val accuracy} & \textbf{Val macro F1} \\
\midrule
1 & 0.6806 & 0.6462 & 0.61 & 0.5400 \\
2 & 0.5769 & 0.6391 & 0.72 & 0.7029 \\
3 & 0.4851 & 0.5709 & 0.79 & 0.7883 \\
4 & 0.3910 & 0.7521 & 0.76 & 0.7520 \\
5 & 0.2880 & 0.5778 & 0.80 & 0.7987 \\
6 & 0.2206 & 0.6792 & 0.80 & 0.7987 \\
7 & 0.1882 & 0.6094 & 0.81 & 0.8091 \\
8 & 0.1462 & 0.6368 & \textbf{0.82} & \textbf{0.8188} \\
\bottomrule
\end{tabular}
\end{center}

Two observations. First, the gap between training loss (0.15) and validation loss (0.64) at epoch 8 shows the model is fitting the small training set closely, yet validation accuracy is still improving rather than degrading, so the run was stopped before overfitting became harmful. Second, the best epoch is the last one, which suggests a longer schedule may still gain a little; we kept 8 epochs as declared in the configuration rather than extending the run after seeing the curve, since tuning the schedule on the basis of these numbers and then reporting them would reintroduce exactly the selection bias discussed below.

\subsection{Test Set Performance}

The selected checkpoint (epoch 8, chosen on the validation split) is evaluated once on the 156-sample held-out test set:

\begin{center}
\begin{tabular}{ll}
\toprule
\textbf{Metric} & \textbf{Value} \\
\midrule
Accuracy                       & \textbf{78.85\%} \\
Macro F1                       & \textbf{0.7880} \\
Macro precision                & 0.7882 \\
Macro recall                   & 0.7879 \\
Evaluation loss (cross-entropy)& 0.7575 \\
Test samples                   & 156 \\
\bottomrule
\end{tabular}
\end{center}

An accuracy of \textbf{78.85\%} with a macro F1 of \textbf{0.7880} on the Bengali test set is a solid result for a low-resource setting, and confirms that IndicBERT representations transfer to Bengali review sentiment without any domain-specific pretraining. Precision and recall are near-identical across the two classes, so the model is not trading one class off against the other.

It is worth stating plainly why this number is lower than the 82.69\% reported in an earlier version of this work. That figure came from a run in which the best epoch was chosen by evaluating on the test set itself, so the test set had already influenced model selection and the resulting score was optimistic. The pipeline described here selects the epoch on a separate validation split and evaluates on the test set exactly once. The 78.85\% is therefore the honest estimate of generalisation, and is the number we stand behind.

\subsection{Additional Metrics}

To provide a more complete evaluation, we report weighted precision, recall, and F1, along with the confusion matrix and full classification report.

\subsubsection{Confusion Matrix}

\begin{center}
\begin{tabular}{lcc}
\toprule
\textbf{True \textbackslash{} Pred} & \textbf{Negative} & \textbf{Positive} \\
\midrule
\textbf{Negative} & 65 & 16 \\
\textbf{Positive} & 17 & 58 \\
\bottomrule
\end{tabular}
\end{center}

\subsubsection{Classification Report}

\begin{lstlisting}
                  precision    recall  f1-score   support

     Negative       0.79      0.80      0.80        81
     Positive       0.78      0.77      0.78        75

     accuracy                           0.79       156
    macro avg       0.79      0.79      0.79       156
weighted avg       0.79      0.79      0.79       156
\end{lstlisting}

\begin{center}
\begin{tabular}{ll}
\toprule
\textbf{Metric} & \textbf{Value} \\
\midrule
Precision (weighted) & \textbf{0.7884} \\
Recall (weighted)    & \textbf{0.7885} \\
F1 (weighted)        & \textbf{0.7884} \\
\bottomrule
\end{tabular}
\end{center}

%-----------------------------------------------------------
\section{Streamlit Application}
%-----------------------------------------------------------

A production-ready Streamlit web application (\texttt{app.py}) is provided alongside the notebook. It offers two modes:

\begin{enumerate}[noitemsep]
    \item \textbf{Single Review Prediction:} The user pastes a Bengali review into a text box and clicks ``Predict''. The app displays the predicted sentiment label and confidence score.
    \item \textbf{Bulk CSV Analysis:} The user uploads a CSV file containing reviews. The app runs inference on all rows, displays a preview table, shows a pie chart of the sentiment distribution, extracts the top keyword themes using TF-IDF n-grams, and provides a downloadable CSV of all predictions.
\end{enumerate}

The model is loaded once and cached using \texttt{@st.cache\_resource} to avoid reloading on each interaction.

To run the app:
\begin{lstlisting}[language=bash]
streamlit run app.py
\end{lstlisting}

\subsection{App Screenshots}

\begin{figure}[H]
    \centering
    \includegraphics[width=0.85\linewidth]{single_review.png}
    \caption{Single review prediction interface.}
\end{figure}

\begin{figure}[H]
    \centering
    \includegraphics[width=0.85\linewidth]{bulk_review.png}
    \caption{Bulk CSV analysis workflow.}
\end{figure}

\begin{figure}[H]
    \centering
    \includegraphics[width=0.7\linewidth]{pie.png}
    \caption{Sentiment distribution pie chart from bulk analysis.}
\end{figure}

%-----------------------------------------------------------
\section{Ablation Study}
%-----------------------------------------------------------

An ablation study was not feasible for this setup. We fine-tuned a single IndicBERT model with the available training configuration, and did not maintain multiple model variants or ablation-friendly checkpoints (e.g., alternative architectures or controlled hyperparameter sweeps). Given the time and resource constraints, we focused on a stable end-to-end pipeline and reporting the primary results.

%-----------------------------------------------------------
\section{Reproducibility and Model Promotion}
%-----------------------------------------------------------

The experiments in this report are produced by a configuration-driven pipeline rather than by an ad-hoc notebook, so that any reported number can be traced back to the exact code, data and hyperparameters that produced it.

\begin{itemize}[noitemsep]
    \item \textbf{Pinned inputs.} The dataset and the base model are both fetched at a fixed Hub commit, so a rerun months later sees the same 998 reviews and the same pre-trained weights.
    \item \textbf{Declared configuration.} All hyperparameters live in \texttt{configs/train.yaml} and any value can be overridden from the command line (\texttt{-o training.epochs=4}). Each run records a hash of its resolved configuration.
    \item \textbf{Tracked runs.} Every run logs its parameters, per-epoch training and validation curves, the git commit, whether the working tree was dirty, and a content fingerprint of each data split to MLflow. The run reported here is \texttt{b85e3fce}, at commit \texttt{3a99a19}.
    \item \textbf{Honest evaluation.} Epoch selection uses the validation split only; the test set is evaluated once, at the end, producing a \texttt{metrics.json} with the metrics, confusion matrix and classification report.
    \item \textbf{Quality gate.} A candidate model is promoted only if it clears absolute floors (accuracy and macro F1 $\geq 0.75$) and does not regress by more than 0.01 macro F1 against the current champion on the same test set. Passing runs are registered and aliased \texttt{@champion}; failing runs exit non-zero and are not published. The gate is what stands between a training run and the model served by the application.
\end{itemize}

This structure is what allowed us to catch the evaluation flaw discussed in Section~\ref{sec:results}: once epoch selection was moved off the test set, the reported accuracy dropped from 82.69\% to 78.85\%, which is the figure that actually reflects unseen data.

%-----------------------------------------------------------
\section{Limitations}
%-----------------------------------------------------------

\begin{enumerate}
    \item \textbf{Small training set:} With only 998 training samples, the model may struggle with rare linguistic phenomena, niche book genres, or very long, complex reviews. Data augmentation or back-translation could help.

    \item \textbf{Domain bias:} The IndicSentiment dataset represents a specific review distribution. The model may perform worse on other Bengali review domains or genres not represented in training.

    \item \textbf{Code-mixed text:} Bengali reviews frequently contain English words written in Bengali script (\textit{Banglish}) or Roman script. The current model is not explicitly trained on code-mixed text and may misclassify such reviews.

    \item \textbf{Binary labels only:} The dataset provides only Positive/Negative labels. Neutral or mixed-sentiment reviews cannot be captured, limiting the nuance of the analyser.

    \item \textbf{No calibration:} The model's raw softmax confidence scores are not calibrated (i.e., a confidence of 0.90 does not necessarily mean the model is correct 90\% of the time). Temperature scaling could be applied post-hoc.

    \item \textbf{Deployment constraints:} The fine-tuned model requires \textasciitilde{}50MB of disk space and a Python environment with PyTorch and HuggingFace Transformers. Running on CPU is possible but slow for bulk inference.
\end{enumerate}

%-----------------------------------------------------------
\section{Conclusion}
%-----------------------------------------------------------

We successfully fine-tuned IndicBERT for Bengali book review sentiment classification, achieving 78.85\% accuracy and 0.788 macro F1 on a held-out test set that was used exactly once, for final evaluation. The approach — leveraging a multilingual, Indic-script-aware pre-trained model and fine-tuning it on task-specific labelled data — is a principled and effective strategy for low-resource NLP. The accompanying Streamlit application makes the model accessible to non-technical users without any coding knowledge. Future work could explore data augmentation, multi-class sentiment (adding a neutral category), and deployment via a REST API for wider integration.

%-----------------------------------------------------------
\begin{thebibliography}{9}

\bibitem{indicsentiment}
Harish Tayyar Madabushi, et al.
\textit{AI4Bharat IndicSentiment Dataset}.
HuggingFace Datasets. 
\url{https://huggingface.co/datasets/ai4bharat/IndicSentiment}

\bibitem{indicbert}
Divyanshu Kakwani, et al.
\textit{IndicNLPSuite: Monolingual Corpora, Evaluation Benchmarks and Pre-trained Multilingual Language Models for Indian Languages}.
EMNLP 2020 Findings.
\url{https://huggingface.co/ai4bharat/indic-bert}

\bibitem{albert}
Zhenzhong Lan, et al.
\textit{ALBERT: A Lite BERT for Self-supervised Learning of Language Representations}.
ICLR 2020.

\bibitem{transformers}
Thomas Wolf, et al.
\textit{HuggingFace Transformers: State-of-the-Art Natural Language Processing}.
EMNLP 2020.

\end{thebibliography}

\end{document}
